% !TeX root = ../main.tex

\ustcsetup{
  keywords  = {加密流量, 流量分类, 自监督学习, 特征融合},
  keywords* = {Encrypted traffic, Traffic classification, Self-Supervised Learning, Feature fusion},
}

\begin{abstract}
  识别流量类型是网络空间安全管理与服务质量控制的核心能力。随着加密技术的普及，以及用户隐私意识的增强，传统基于端口和深度包检测（DPI）的网络流量分类方法因载荷不可见而失效。基于机器学习方法虽利用数据流的统计特征实现分类，但却嫉妒依赖人工特征工程，泛化能力有限且实施成本高。基于深度学习技术通过端到端学习范式突破上述限制，可直接处理原始流量数据，自动提取特征，并减少对专家知识的依赖。但现有方案多存在特征表征不足、易受网络环境波动干扰、过度依赖隐私载荷以及特征维度单一等局限。针对上述问题，本文的研究工作主要基于两个阶段的核心工作展开。第一阶段针对单一维度特征提取的不稳定性，研究基于自监督学习的流级别鲁棒特征提取技术；第二阶段针对隐私合规以及数据特征维度单一问题，构建了仅依赖数据报头的包级别特征提取技术以及双模态特征融合分类框架。通过这两个阶段的递进式研究，本文提出了一种融合流级别时序特征与包级别结构特征的深度学习加密流量分类方法。具体研究内容如下：

  （1）在本文中，我们首先提出了基于 BYOL 自监督架构的流级别鲁棒特征分类方法。针对网络环境波动导致的流量时序特征失真，设计了噪声注入、尺度缩放与随机遮蔽两种增强函数，通过非对称网络构建动态系统，从而提取出两种增强变体之间不变的，具有鲁棒性的流级别特征。实验结果表明，在无需负样本的情况下，该方法有效克服了表征崩溃问题，与不具有对比学习机制的传统深度学习方法相比，在多个公开加密流量数据集上分类准确率提升了约 5.2\%，显著增强了模型在复杂网络环境下的鲁棒性。

  （2）为了进一步提升分类效果，且更好地缓解单模态特征分类存在的类别混淆问题，本文构建了兼顾隐私保护与多维表征的双模态融合分类框架。在包级别处理中，仅利用 IP 数据报头信息结合 n-gram 嵌入，实现在规避敏感载荷的同时捕捉协议微观特征，使得模型有效地学习了报头字段的字节组合与顺序特征。此外还设计了基于交叉注意力机制的融合模块，实现流级宏观时序特征与包级微观结构特征的深度语义对齐。实验证明，该双模态融合模型在公开标准数据集上的准确率均达到 94\% 以上，较单一模态方法提升了 5.8\%-9.4\%，相比现有主流双模态方法在综合性能上优化了约 4.1\%。  

\end{abstract}

\begin{abstract*}
  Identifying traffic types is a fundamental and crucial capability in network research, essential for applications such as quality of service control and security testing. With increasing user privacy awareness and the widespread adoption of encryption tools like SSL/TLS, VPNs, and Tor, traditional port-based and deep packet inspection (DPI) methods have become ineffective due to their inability to parse encrypted packet payloads. While machine learning methods utilize flow statistics for classification, they rely on manual feature engineering, resulting in limited generalization capabilities and high implementation costs. Deep learning technology overcomes these limitations through an end-to-end learning paradigm: it can directly process raw traffic data, automatically extract features, support early classification, and reduce reliance on expert knowledge. However, existing deep learning solutions often have significant limitations: they analyze and identify traffic based solely on flow-level or packet-level features, neglecting traffic features contained in another dimension, failing to integrate dual-dimensional information; furthermore, in packet-level processing, there are instances of mixed use of packet headers and payloads, failing to distinguish functional differences and infringing on user privacy by using payload data. Therefore, researching a deep learning-based fusion encrypted traffic classification method that integrates flow-level and packet-level features, using only IP packet headers and excluding packet payloads, is of significant value for improving network security and service quality. The main contributions of this paper are as follows:

  (1) In this paper, we first propose a flow-level feature classification method based on the BYOL model. Our goal is to achieve effective early traffic classification using only the early packet length sequence of the data flow. This scheme extracts robust flow-level features invariant between the two augmentation variants by truncating the packet length sequence and then applying two differential augmentations to the same length sequence, respectively, and inputting these into the online and target networks of the BYOL model. Experiments show that our method achieves higher classification accuracy on multiple encrypted traffic datasets, improving it by approximately $5\%$ compared to traditional deep learning methods without contrastive learning mechanisms.

  (2) To further improve classification performance and better alleviate the category confusion problem in single-modal feature classification, this paper designs a fusion-based encrypted traffic classification method. In addition to flow-level features, we further extract packet-level features. By using only IP datagram header information and applying n-gram embeddings with different window sizes, the model effectively learns the byte combination and order features of the header fields. Furthermore, through an attention-based feature fusion module, flow-level and packet-level features are adaptively weighted and fused, effectively modeling the non-linear relationships between features. The method is validated on multiple encrypted traffic datasets. Experimental results show that the classification method fusing flow-level and packet-level features significantly outperforms methods using only single-dimensional features, improving the accuracy of encrypted traffic classification by approximately $5\%-9\%$. Compared to current representative bimodal methods, it improves accuracy by approximately $4\%$.
\end{abstract*}
